<<<<<<< HEAD
\section{Results}

This section presents the experimental validation of the 3-user co-located VR system during an extended co-located run. The longest continuous logging window with all three headsets active lasted 85.1 minutes (5073 samples per headset at 1Hz; 15,234 total samples across devices). The host headset additionally logged 97.0 minutes including pre-session setup and post-session. Results are organized by research question.

\subsection{RQ1, RQ3 \& RQ4: Co-location Implementation, Monitoring, and Deployment}

Implementation details for colocation, logging, and deployment are described in Sections~\ref{sec:prototype-colocation}, \ref{sec:prototype-logging}, and \ref{sec:prototype-deployment}. In summary, the prototype repeatedly established a shared coordinate system for three headsets using a host-advertised local session (group UUID) and a distributed shared spatial anchor, enabling consistent placement of the \texttt{Arena} and shared content. Continuous monitoring was achieved through on-device 1Hz logging (CSV + JSON metadata) across the full session duration. For reproducible multi-headset iteration, APK installation and launch were standardized using Meta Quest Developer Hub (MQDH) with Developer Mode enabled on each headset.

\subsection{RQ2: Calibration Accuracy and Drift}

Automated spatial anchor calibration achieved high accuracy under stationary baseline conditions. During the 85.1-minute 3-headset run, the client headset that logged non-zero calibration error reported a mean error of $1.21 \pm 0.90$mm with a maximum of 5.78mm. This remains well within Reimer et al.'s~\cite{ReimerDennis2021CfSV} $<10$mm safety-critical reference threshold for collision avoidance, supporting the feasibility of anchor-based co-location.

Here, \emph{non-zero calibration error} means that the headset recorded a non-zero \emph{residual alignment error} in our logs (i.e., a non-zero Euclidean distance between the headset's local anchor position and the shared anchor position after alignment). In this run, only one headset recorded calibration error values above zero in the CSV logs.

It should be noted that this drift evaluation was conducted with headsets placed stationary on a table. Consequently, these results represent a best-case baseline and may underestimate drift and transient misalignment that can occur under active user motion and re-localization events.

\begin{table}[h]
\centering
\caption{Calibration Accuracy (Headset with Non-Zero Logging)}
\begin{tabular}{|l|c|c|c|}
\hline
Metric & \textbf{Mean (mm)} & \textbf{SD} & \textbf{Max} \\
\hline
Error & 1.21 & 0.90 & 5.78 \\
\hline
\end{tabular}
\end{table}

Under these stationary conditions, calibration drift remained low over the full 85.1-minute 3-headset run, with maximum error staying well below the 10mm reference threshold.

\subsection{RQ5: Feasibility Benchmarks and Long-Term Stability}

\subsubsection{Network Latency Performance}

Network latency measurements show that the system's mean RTT remained below Van Damme et al.'s~\cite{VanDammeSam2024Iolo} target threshold for good collaborative quality of experience (QoE $\leq$ 75ms RTT). Over the 85.1-minute 3-headset run, mean RTT was $58.2 \pm 21.3$ms with a 95th percentile of 86.0ms. While the mean remains within the 75ms target, the tail latency exceeds it at P95 and includes occasional spikes (max 791.4ms) that occurred during initialization and synchronization events.
=======
% \section{Results}
% Describe the outcome of your experiment. Discuss the success and failure of your experiments.
\section{Results}

This section presents the experimental validation of the 3-user co-located VR system during an extended co-located run. The longest continuous logging window with all three headsets active lasted 85.1 minutes (5073 samples per headset at 1Hz; 15,234 total samples across devices). The host headset additionally logged 97.0 minutes including pre-session setup and post-session tail time. Results are organized by research question.

\subsection{RQ1: Technical Benchmark Achievement}

\subsubsection{Network Latency Performance}

Network latency measurements show that the system's \emph{mean} RTT remained below Van Damme et al.'s~\cite{VanDammeSam2024Iolo} target threshold for good collaborative quality of experience (QoE $\leq$ 75ms RTT). Over the 85.1-minute 3-headset run, mean RTT was $58.2 \pm 21.3$ms with a 95th percentile of 86.0ms. While the mean remains within the 75ms target, the tail latency exceeds it at P95 and includes occasional spikes (max 791.4ms) during initialization and synchronization events.
>>>>>>> 32df077fa8bbb82b5374f4992fe0114abd8167f7

\begin{table}[h]
\centering
\caption{Network Latency Performance}
\begin{tabular}{|l|c|c|c|}
\hline
<<<<<<< HEAD
                    	\textbf{Metric} & \textbf{Mean (ms)} & \textbf{SD} & \textbf{P95 (ms)} \\
=======
			\textbf{Metric} & \textbf{Mean (ms)} & \textbf{SD} & \textbf{P95 (ms)} \\
>>>>>>> 32df077fa8bbb82b5374f4992fe0114abd8167f7
\hline
Overall & 58.2 & 21.3 & 86.0 \\
\hline
\end{tabular}
\end{table}

<<<<<<< HEAD
Despite occasional outliers, the system consistently maintained performance below the 75ms threshold, supporting the feasibility of consumer wireless networking for co-located collaborative training scenarios when supported by robust error handling.

\subsubsection{Frame Rate Stability}

Frame rate performance was evaluated against the Meta Quest 3's 72Hz refresh rate target set for this study. The system achieved a mean frame rate of $72.2 \pm 0.7$fps. 
=======
Despite occasional outliers, the system consistently maintained performance below the 75ms threshold, validating consumer wireless networking for safety-critical training applications when supported by robust error handling.

\subsubsection{Frame Rate Stability}

Frame rate performance was evaluated against the Meta Quest 3's native 72Hz refresh rate target. The system achieved a mean frame rate of $72.2 \pm 0.7$fps, remaining tightly clustered around the 72Hz target throughout the run.
>>>>>>> 32df077fa8bbb82b5374f4992fe0114abd8167f7

\begin{table}[h]
\centering
\caption{Frame Rate Performance Summary}
\begin{tabular}{|l|c|c|}
\hline
\textbf{State} & \textbf{Frame Rate (fps)} & \textbf{Notes} \\
\hline
Overall Mean & $72.2 \pm 0.7$ & Stable Performance \\
5th--95th Percentile & $70.9 - 73.1$ & Includes initialization variance \\
\hline
\end{tabular}
\end{table}

<<<<<<< HEAD
Frame rates remained tightly clustered around the 72Hz target during the run, indicating stable performance under the tested load.


=======
\subsubsection{Calibration Accuracy and Drift}

Automated spatial anchor calibration achieved high accuracy under stationary baseline conditions. During the 85.1-minute 3-headset run, the client headset that logged non-zero calibration error reported a mean error of $1.21 \pm 0.90$mm with a maximum of 5.78mm. This remains within Reimer et al.'s~\cite{ReimerDennis2021CfSV} $<10$mm safety-critical reference threshold, supporting the feasibility of anchor-based co-location.

\begin{table}[h]
\centering
\caption{Calibration Accuracy (Headset with Non-Zero Logging)}
\begin{tabular}{|l|c|c|c|}
\hline
\textbf{Metric} & \textbf{Mean (mm)} & \textbf{SD} & \textbf{Max} \\
\hline
Error & 1.21 & 0.90 & 5.78 \\
\hline
\end{tabular}
\end{table}

Calibration drift was minimal over the full 85.1-minute 3-headset run, with maximum error remaining well below the 10mm reference threshold.

\subsection{RQ2 \& RQ3: Collaboration Effectiveness}

\subsubsection{Task Performance}

While specific WIM interface metrics were not collected in this technical validation phase, the system successfully supported 3-user co-located interaction for an 85.1-minute run without technical failure. Under stationary baseline conditions, the shared coordinate system remained within single-digit millimeter error on the headset reporting non-zero calibration error.

\subsection{RQ4: Long-Term Performance Stability}
>>>>>>> 32df077fa8bbb82b5374f4992fe0114abd8167f7

\subsubsection{Thermal Performance}

Device temperature monitoring revealed consistent thermal buildup patterns. Headsets reached a maximum temperature of $52.6^{\circ}$C.

<<<<<<< HEAD
Temperature correlated weakly with performance (frame rate vs. temperature: $r=0.14$).

Despite significant thermal buildup, the system maintained stable performance (72fps and calibration error remaining well below the 10mm threshold) throughout the extended sessions, demonstrating the Quest 3's effective thermal management for long-duration training.
=======
Temperature correlated weakly with performance metrics:
\begin{itemize}
    \item Frame rate vs Temp: $r=0.14$ (Weak positive correlation)
\end{itemize}

Despite significant thermal buildup, the system maintained stable performance throughout the extended run, demonstrating the Quest 3's effective thermal management for long-duration training.
>>>>>>> 32df077fa8bbb82b5374f4992fe0114abd8167f7
